% ============================================================================
% MUSTERLOSUNG: Stundenleistung Mi casa
% ============================================================================
% Sequenz 3, Block d: Musterlosung zur Stundenleistung
% ============================================================================

\documentclass[11pt, a4paper]{article}

% ============================================================================
% SW-MODUS TOGGLE
% ============================================================================
\newif\ifswmodus
\swmodusfalse

% ============================================================================
% PAKETE
% ============================================================================

\usepackage[utf8]{inputenc}
\usepackage[T1]{fontenc}
\usepackage[ngerman]{babel}
\usepackage{helvet}
\renewcommand{\familydefault}{\sfdefault}

\usepackage[margin=2cm]{geometry}
\usepackage{enumitem}
\usepackage{tabularx}
\usepackage{booktabs}
\usepackage{multirow}
\usepackage{pgffor}

\usepackage{xcolor}
\usepackage{framed}
\usepackage{tcolorbox}
\tcbuselibrary{skins}

\usepackage{fancyhdr}
\usepackage{lastpage}
\usepackage{needspace}

% ============================================================================
% FARBEN
% ============================================================================

\definecolor{spanisch-color}{HTML}{c41e3a}
\definecolor{grau-color}{HTML}{888888}
\definecolor{hellgrau-color}{HTML}{f5f5f5}
\definecolor{orange-color}{HTML}{b35c00}
\definecolor{loesung-color}{HTML}{c62828}

\definecolor{sw-dark}{HTML}{000000}
\definecolor{sw-gray}{HTML}{666666}
\definecolor{sw-white}{HTML}{ffffff}

\ifswmodus
    \colorlet{spanisch}{sw-dark}
    \colorlet{grau}{sw-gray}
    \colorlet{hellgrau}{sw-white}
    \colorlet{orange}{sw-dark}
    \colorlet{loesung}{sw-dark}
\else
    \colorlet{spanisch}{spanisch-color}
    \colorlet{grau}{grau-color}
    \colorlet{hellgrau}{hellgrau-color}
    \colorlet{orange}{orange-color}
    \colorlet{loesung}{loesung-color}
\fi

\newcommand{\fachfarbe}{spanisch}

% ============================================================================
% VARIABLEN
% ============================================================================

\newcommand{\thema}{Mi casa -- Stundenleistung}
\newcommand{\sequenz}{03}
\newcommand{\block}{d}
\newcommand{\fach}{Spanisch}
\newcommand{\klasse}{7}
\newcommand{\maxpunkte}{20}
\newcommand{\dauer}{25 Minuten}
\newcommand{\datum}{\today}

% ============================================================================
% KOPF- UND FUSSZEILE
% ============================================================================

\pagestyle{fancy}
\fancyhf{}
\renewcommand{\headrulewidth}{0.4pt}
\renewcommand{\footrulewidth}{0.4pt}

\lhead{\fach{} -- Klasse \klasse}
\chead{\textcolor{loesung}{\textbf{SL-\sequenz\block{} -- MUSTERLOSUNG}}}
\rhead{\datum}

\lfoot{}
\cfoot{}
\rfoot{Seite \thepage\ von \pageref{LastPage}}

% ============================================================================
% BEFEHLE
% ============================================================================

\newcommand{\punktebox}[1]{\hfill\fbox{\small \_/#1}}
\newcommand{\loesung}[1]{\textcolor{loesung}{\textbf{#1}}}
\newcommand{\luecke}[1]{\rule{#1}{0.4pt}}

\newtcolorbox{vokabelkasten}{
    colback=white,
    colframe=\fachfarbe,
    boxrule=0.5pt,
    arc=0pt,
    left=5pt, right=5pt, top=5pt, bottom=5pt
}

\newenvironment{aufgabe}[1]{%
    \needspace{5\baselineskip}%
    \nopagebreak[4]%
    \vspace{10pt}
    \noindent\textbf{\textcolor{\fachfarbe}{Aufgabe #1}}
    \nopagebreak[4]%
    \vspace{5pt}
    \nopagebreak[4]%
    \begin{list}{}{\leftmargin=0pt}
    \item[]
}{%
    \end{list}
}

\newcommand{\es}[1]{\textcolor{\fachfarbe}{\textbf{#1}}}

% ============================================================================
% DOKUMENT
% ============================================================================

\begin{document}

% ============================================================================
% KOPFBEREICH
% ============================================================================

\begin{center}
    {\Large\bfseries\textcolor{\fachfarbe}{Stundenleistung: \thema}}\\[0.3em]
    {\large\textcolor{loesung}{\textbf{--- MUSTERLOSUNG ---}}}
\end{center}

\vspace{5pt}

\noindent
\begin{tabularx}{\textwidth}{|X|X|X|X|}
\hline
\textbf{Name:} \textcolor{loesung}{Musterschuler/in} &
\textbf{Klasse:} \klasse &
\textbf{Datum:} \datum &
\textbf{Zeit:} \dauer \\
\hline
\end{tabularx}

\vspace{15pt}

% ============================================================================
% AUFGABE 1: Raume zuordnen (4 Punkte)
% ============================================================================

\begin{aufgabe}{1: Raume zuordnen}
Ordne die spanischen Raume den deutschen Ubersetzungen zu. Schreibe den passenden Buchstaben in die Lucke. \punktebox{4}
\end{aufgabe}

\begin{center}
\begin{tabular}{lcp{6cm}}
1. \es{la cocina} & = \loesung{c} & a) das Schlafzimmer \\[0.8em]
2. \es{el dormitorio} & = \loesung{a} & b) das Badezimmer \\[0.8em]
3. \es{el salon} & = \loesung{d} & c) die Kuche \\[0.8em]
4. \es{el bano} & = \loesung{b} & d) das Wohnzimmer \\
\end{tabular}
\end{center}

\textit{\textcolor{grau}{Bewertung: 1 Punkt pro korrekter Zuordnung = 4 Punkte}}

\vspace{15pt}

% ============================================================================
% AUFGABE 2: hay + Mobel (4 Punkte)
% ============================================================================

\begin{aufgabe}{2: hay + Mobel}
Erganze die Satze mit \es{hay} und einem passenden Mobel aus dem Kasten. \punktebox{4}
\end{aufgabe}

\begin{vokabelkasten}
\textit{Mobel:} una cama | un sofa | una mesa | un armario
\end{vokabelkasten}

\vspace{0.5em}

\noindent
a) En la cocina \loesung{hay} \loesung{una mesa}. \\[1em]
b) En el dormitorio \loesung{hay} \loesung{una cama}. \\[1em]
c) En el salon \loesung{hay} \loesung{un sofa}. \\[1em]
d) En el dormitorio tambien \loesung{hay} \loesung{un armario}. \\

\textit{\textcolor{grau}{Bewertung: 0,5P fur hay + 0,5P fur passendes Mobel = 4 Punkte}}

\vspace{15pt}

% ============================================================================
% AUFGABE 3: Farben + Adjektivanpassung (6 Punkte)
% ============================================================================

\begin{aufgabe}{3: Farben + Adjektivanpassung}
Erganze die Farbadjektive. Achte auf die richtige Endung! \punktebox{6}
\end{aufgabe}

\begin{vokabelkasten}
\textit{Farben:} verde | rojo/roja | blanco/blanca | marron/marrones | azul | grande
\end{vokabelkasten}

\vspace{0.5em}

\noindent
a) El sofa es \loesung{verde} (grun). \\[1em]
b) La silla es \loesung{roja} (rot). \\[1em]
c) Las mesas son \loesung{blancas} (weiß). \\[1em]
d) Los armarios son \loesung{marrones} (braun). \\[1em]
e) La cama es \loesung{azul} (blau). \\[1em]
f) El jardin es \loesung{grande} (groß). \\

\textit{\textcolor{grau}{Bewertung: 1 Punkt pro korrektem Adjektiv (inkl. Endung) = 6 Punkte}}

\textbf{Hinweise zur Korrektur:}
\begin{itemize}
    \item a) \textit{verde} -- unveranderlich (Konsonant-Endung)
    \item b) \textit{roja} -- feminin, -a Endung
    \item c) \textit{blancas} -- feminin Plural
    \item d) \textit{marrones} -- maskulin Plural (-on $\rightarrow$ -ones)
    \item e) \textit{azul} -- unveranderlich
    \item f) \textit{grande} -- unveranderlich (vor Nomen: gran)
\end{itemize}

\vspace{15pt}

% ============================================================================
% AUFGABE 4: Zimmer beschreiben (6 Punkte)
% ============================================================================

\begin{aufgabe}{4: Zimmer beschreiben}
Beschreibe ein Zimmer deiner Wahl auf Spanisch. Schreibe mindestens 4 Satze. Verwende \es{hay}, Mobel und Farben. \punktebox{6}
\end{aufgabe}

\noindent\textbf{Mi habitacion favorita:} \loesung{el dormitorio}

\textbf{Beispiellosung:}

\loesung{En mi dormitorio hay una cama grande. La cama es blanca.}

\loesung{Tambien hay un armario marron. El armario es muy grande.}

\loesung{Hay una mesa pequena. La mesa es negra.}

\loesung{En mi dormitorio hay una silla azul. Me gusta mi dormitorio.}

\vspace{1em}

\textbf{Bewertungskriterien:}

\begin{tabular}{|l|c|p{8cm}|}
\hline
\textbf{Kriterium} & \textbf{Punkte} & \textbf{Beschreibung} \\
\hline
Mindestens 4 Satze & 2P & 0,5P pro Satz bis max. 2P \\
\hline
Korrektes \textit{hay} & 2P & Mind. 2x korrekt verwendet \\
\hline
Farbadjektive & 1P & Mind. 2 Farben, korrekt angepasst \\
\hline
Sprachliche Korrektheit & 1P & Genus, Artikel, Satzstellung \\
\hline
\end{tabular}

\vspace{1em}

\textbf{Abzuge:}
\begin{itemize}
    \item Fehlender Satz: -0,5P (unter 4 Satze)
    \item Falsches Genus bei Mobeln: -0,5P
    \item Adjektiv nicht angepasst: -0,5P
    \item hay mit falschem Artikel (*hay la): -0,5P
\end{itemize}

% ============================================================================
% PUNKTEUBERSICHT
% ============================================================================

\vspace{20pt}
\hrule
\vspace{10pt}

\noindent
\begin{tabularx}{\textwidth}{|X|c|c|c|c||c|}
\hline
& \textbf{Aufg. 1} & \textbf{Aufg. 2} & \textbf{Aufg. 3} & \textbf{Aufg. 4} & \textbf{Gesamt} \\
\hline
\textbf{Punkte} & \loesung{4}/4 & \loesung{4}/4 & \loesung{6}/6 & \loesung{6}/6 & \loesung{20}/\maxpunkte \\
\hline
\end{tabularx}

\vspace{15pt}

% ============================================================================
% BEWERTUNGSSKALA
% ============================================================================

\section*{Bewertungsskala (14-NP, Sek I)}

\begin{center}
\begin{tabular}{|c|c|c|c|c|c|c|c|c|c|c|c|c|c|c|}
\hline
\textbf{NP} & 14 & 13 & 12 & 11 & 10 & 9 & 8 & 7 & 6 & 5 & 4 & 3 & 2 & 1 \\
\hline
\textbf{\%} & 100 & 96 & 91 & 86 & 80 & 73 & 67 & 60 & 53 & 47 & 40 & 33 & 27 & 20 \\
\hline
\textbf{BE} & 20 & 19 & 18 & 17 & 16 & 15 & 13 & 12 & 11 & 9 & 8 & 7 & 5 & 4 \\
\hline
\end{tabular}
\end{center}

\vspace{0.5em}

\textbf{Umrechnung in Schulnoten:}
\begin{center}
\begin{tabular}{|c|c|c|c|c|c|c|}
\hline
\textbf{Note} & 1 & 2 & 3 & 4 & 5 & 6 \\
\hline
\textbf{NP} & 13-14 & 10-12 & 7-9 & 4-6 & 1-3 & 0 \\
\hline
\textbf{BE} & 19-20 & 15-18 & 11-14 & 7-10 & 4-6 & 0-3 \\
\hline
\end{tabular}
\end{center}

\end{document}
